\documentclass[10pt,letterpaper,twocolumn]{article}

% Paquetes: se eliminaron duplicados y se compactó el formato
\usepackage[spanish,es-nodecimaldot]{babel}
\usepackage[utf8]{inputenc}
\usepackage[T1]{fontenc}
\usepackage{lmodern,microtype}
\usepackage{graphicx}
\usepackage{float}
\usepackage{subcaption}
\usepackage{amsmath,amssymb,booktabs,tabularx,array}
\usepackage{geometry,caption,siunitx,hyperref}
\usepackage{algorithm,algpseudocode}
\usepackage{titlesec}
\usepackage{balance}
\usepackage{etoolbox}

\geometry{letterpaper,margin=1.45cm}
\setlength{\columnsep}{0.55cm}
\setlength{\parindent}{1em}
\setlength{\parskip}{1pt}
\linespread{0.97}
\raggedbottom

% Espaciado compacto para conservar el límite de páginas
\titlespacing*{\section}{0pt}{0.9ex plus 0.2ex minus 0.2ex}{0.45ex}
\titlespacing*{\subsection}{0pt}{0.7ex plus 0.2ex minus 0.2ex}{0.3ex}
\setlength{\textfloatsep}{5pt plus 1pt minus 1pt}
\setlength{\floatsep}{4pt plus 1pt minus 1pt}
\setlength{\intextsep}{4pt plus 1pt minus 1pt}
\setlength{\abovedisplayskip}{3pt plus 1pt minus 1pt}
\setlength{\belowdisplayskip}{3pt plus 1pt minus 1pt}
\setlength{\abovedisplayshortskip}{2pt plus 1pt minus 1pt}
\setlength{\belowdisplayshortskip}{2pt plus 1pt minus 1pt}
\algrenewcommand\algorithmicrequire{\textbf{Entrada:}}
\algrenewcommand\algorithmicensure{\textbf{Salida:}}
\floatname{algorithm}{Algoritmo}
\captionsetup{font=footnotesize,labelfont=bf,skip=2pt}\addto\captionsspanish{\renewcommand{\tablename}{Tabla}}
\captionsetup[subfigure]{font=scriptsize,skip=1pt}
\sisetup{output-decimal-marker={.}}
\hypersetup{hidelinks}
\AtBeginEnvironment{algorithm}{\footnotesize}

% Si las imágenes todavía no están en la carpeta, se muestra un recuadro temporal.
% Al colocar los PNG con esos nombres, se insertarán automáticamente.
\newcommand{\grafica}[1]{%
  \IfFileExists{#1}{\includegraphics[width=\columnwidth,height=4.1cm,keepaspectratio]{#1}}%
  {\fbox{\parbox[c][4.1cm][c]{0.94\columnwidth}{\centering\scriptsize Gráfica pendiente: \texttt{\detokenize{#1}}}}}%
}

\begin{document}

% Encabezado a una columna
\twocolumn[{
\begin{center}
{\Large\bfseries CLASIFICACIÓN DEL ESTADO FETAL MEDIANTE CARDIOTOCOGRAFÍA Y MACHINE LEARNING\par}
\vspace{0.35em}
{\small Dylan Luisjuan Rico \quad Leonardo Alejandro Ambriz Mendoza\\
Derek Lugo Aguirre \quad Ricardo Chávez Altamirano\par}
\vspace{0.2em}
{\footnotesize Centro Universitario de Guadalajara (CUGDL), Universidad de Guadalajara\par}
\end{center}
\vspace{0.2em}
\begin{minipage}{\textwidth}
\centering\textbf{Resumen}\par
\small
Este trabajo compara algoritmos de aprendizaje automático para clasificar registros cardiotocográficos (CTG) en tres categorías: normal, sospechoso y patológico. Se implementaron K vecinos más cercanos (KNN) y regresión logística multiclase uno contra el resto (OvR), tanto mediante código propio como con \texttt{scikit-learn}. Se describe el preprocesamiento, la evaluación predictiva y la medición de tiempos de entrenamiento y predicción. También se analiza la complejidad temporal teórica y su relación con los resultados experimentales. KNN obtuvo el mejor desempeño global de clasificación, mientras que \texttt{scikit-learn} redujo notablemente el tiempo de predicción de KNN. Los resultados corresponden a una evaluación experimental y no constituyen un sistema de diagnóstico clínico autónomo.\par
\vspace{0.2em}
\noindent\textbf{Palabras clave:} aprendizaje automático, cardiotocografía, KNN, regresión logística, clasificación, complejidad temporal.
\end{minipage}
\vspace{0.5em}
}]

\section{Introducción}
La cardiotocografía (CTG) registra la frecuencia cardiaca fetal y la actividad uterina. Su interpretación permite identificar patrones asociados con estados normales o posibles señales de compromiso fetal (Campos \& Bernardes, 2000). Debido a la cantidad de características y a la variabilidad de interpretación, el aprendizaje automático puede servir como apoyo para clasificar estos registros. El objetivo de este proyecto es comparar KNN y regresión logística multiclase OvR, mediante implementaciones propias y de \texttt{scikit-learn}, evaluando tanto su desempeño predictivo como sus tiempos de entrenamiento y predicción.

\section{Planteamiento del problema}
Se busca clasificar cada registro CTG en una de tres categorías: normal, sospechoso o patológico, a partir de sus características numéricas. El conjunto contiene 2\,126 registros y 21 variables predictoras. La distribución es desigual: 77.85\% normal, 13.88\% sospechoso y 8.28\% patológico (Campos \& Bernardes, 2000). Por ello, además de la exactitud, se consideran F1 macro y Recall por clase, ya que una exactitud elevada puede ocultar un desempeño limitado en las clases minoritarias.

\section{Metodología}
\subsection{Datos y preparación}
Se utilizó \texttt{CTG\_clean.csv}, con 2\,126 observaciones, 21 características y una variable objetivo NSP. Para la evaluación predictiva se separó 80\% para entrenamiento y 20\% para prueba, con semilla 42: 1\,701 registros de entrenamiento y 425 de prueba. Las características se estandarizaron mediante
\[
z=\frac{x-\mu_{\mathrm{train}}}{\sigma_{\mathrm{train}}},
\]
calculando media y desviación estándar solo con entrenamiento para evitar fuga de información (Pedregosa et al., 2011).

El experimento temporal fue independiente: se evaluaron diez fracciones del conjunto (10\% a 100\%) mediante particiones estratificadas 80/20. En el tamaño del 100\% se utilizaron 1\,700 registros para entrenamiento y 426 para prueba; estos tamaños no deben confundirse con los de la evaluación predictiva.

\subsection{K vecinos más cercanos (KNN)}
KNN asigna a cada observación la clase mayoritaria entre sus $k$ vecinos más cercanos. Se utilizó $k=5$, distancia euclidiana y votación uniforme (Scikit-learn Developers, 2026b):
\[
d(x,x_i)=\sqrt{\sum_{j=1}^{p}(x_j-x_{ij})^2}.
\]

\begin{algorithm}[H]
\caption{Clasificador KNN propio}
\label{alg:knn}
\begin{algorithmic}[1]
\Require $X_{train}$, $y_{train}$, número de vecinos $k=5$
\Ensure Etiquetas predichas para $X$
\Function{Fit}{$X,y$}
  \State $X_{train}\gets X$; $y_{train}\gets y$
  \State \Return modelo $(X_{train},y_{train})$
\EndFunction
\Function{DistanciaEuclidiana}{$a,b$}
  \State \Return $\sqrt{\sum_{j=1}^{p}(a_j-b_j)^2}$
\EndFunction
\Function{PredecirUno}{$x$}
  \State $D\gets[\,]$
  \ForAll{$x_i\in X_{train}$}
    \State Agregar $\Call{DistanciaEuclidiana}{x,x_i}$ a $D$
  \EndFor
  \State $I\gets\Call{OrdenarIndices}{D}$; $I_k\gets$ primeros $k$ índices
  \State $Y_k\gets y_{train}[I_k]$
  \State \Return $\Call{Moda}{Y_k}$
\EndFunction
\Function{Predict}{$X$}
  \State $P\gets[\,]$
  \ForAll{$x\in X$}
    \State Agregar $\Call{PredecirUno}{x}$ a $P$
  \EndFor
  \State \Return arreglo $P$
\EndFunction
\end{algorithmic}
\end{algorithm}

En la implementación propia, \texttt{fit} almacena referencias a los datos; bajo esta condición, su costo adicional es $O(1)$. Para cada una de las $M$ observaciones de prueba se calculan $N$ distancias de costo $O(p)$ y se ordenan los $N$ índices, por lo que:
\[
T_{\mathrm{predict}}=O\!\left(M(Np+N\log N)\right).
\]
Con $p$ fijo, el término dominante es $O(MN\log N)$. La versión de \texttt{scikit-learn} utilizó distancia euclidiana, \texttt{algorithm=auto} y \texttt{n\_jobs=-1}. Con 21 características se considera búsqueda por fuerza bruta, cuyo costo de predicción es aproximadamente $O(MNp)$. La vectorización y el paralelismo pueden reducir el tiempo observado, pero no eliminan el trabajo asintótico. El ajuste de la biblioteca incluye validación y preparación de entradas, por lo que su costo puede aproximarse por $O(Np)$ (Scikit-learn Developers, 2026b).

\subsection{Regresión logística multiclase OvR}
La regresión logística se extendió a tres clases mediante One-vs-Rest: se entrena un clasificador binario por clase y se asigna la clase con mayor probabilidad estimada (Scikit-learn Developers, 2026a). La implementación propia empleó tasa de aprendizaje $\alpha=0.1$ y 2\,000 iteraciones de descenso de gradiente por lotes. El costo auxiliar corresponde a entropía cruzada binaria; las probabilidades se limitan a $[\epsilon,1-\epsilon]$ para evitar logaritmos de cero.

\begin{algorithm}[H]
\caption{Regresión logística binaria propia}
\label{alg:logistic}
\begin{algorithmic}[1]
\Require $X$, $y$, tasa $\alpha=0.1$, iteraciones $T=2000$
\Ensure Pesos $w$ y sesgo $b$ entrenados
\Function{Probabilidad}{$X,w,b$}
  \State \Return $\sigma(Xw+b)$, donde $\sigma(z)=1/(1+e^{-z})$
\EndFunction
\Function{Costo}{$X,y,w,b$}
  \State $\epsilon\gets10^{-15}$; $p\gets\operatorname{clip}(\Call{Probabilidad}{X,w,b},\epsilon,1-\epsilon)$
  \State \Return $-\frac{1}{N}\sum_{i=1}^{N}[y_i\log p_i+(1-y_i)\log(1-p_i)]$
\EndFunction
\Function{Gradientes}{$X,y,w,b$}
  \State $e\gets\Call{Probabilidad}{X,w,b}-y$; $N\gets$ filas de $X$
  \State \Return $X^Te/N,\;\sum e/N$
\EndFunction
\Function{Fit}{$X,y$}
  \State $w\gets\mathbf{0}_{p}$; $b\gets0$
  \For{$t=1$ \textbf{hasta} $T$}
    \State $(dw,db)\gets\Call{Gradientes}{X,y,w,b}$
    \State $w\gets w-\alpha dw$; $b\gets b-\alpha db$
  \EndFor
  \State \Return $(w,b)$
\EndFunction
\Function{Predict}{$X,w,b$}
  \State \Return $(\Call{Probabilidad}{X,w,b}\geq0.3)$ convertido a enteros
\EndFunction
\end{algorithmic}
\end{algorithm}

\begin{algorithm}[H]
\caption{Clasificación multiclase mediante OvR}
\label{alg:ovr}
\begin{algorithmic}[1]
\Require $X$, etiquetas $y$, tasa $\alpha$, iteraciones $T$
\Ensure Predicciones multiclase
\Function{FitOvR}{$X,y$}
  \State $C\gets$ valores únicos de $y$; $M\gets[\,]$
  \ForAll{$c\in C$}
    \State $y_c\gets(y=c)$ convertido a enteros
    \State Entrenar modelo $m_c$ mediante Algoritmo~\ref{alg:logistic}
    \State Agregar $m_c$ a $M$
  \EndFor
  \State \Return $(C,M)$
\EndFunction
\Function{PredictOvR}{$X,C,M$}
  \For{$j=1$ \textbf{hasta} $|C|$}
    \State $P_{:,j}\gets$ probabilidades del modelo $M_j$ para $X$
  \EndFor
  \State \Return $C[\operatorname{argmax}_{j}(P_{i,j})]$ para cada fila $i$
\EndFunction
\end{algorithmic}
\end{algorithm}

Para la implementación propia, el entrenamiento de los $C$ clasificadores requiere $O(CINp)$, donde $I$ es el número de iteraciones, $N$ las muestras y $p$ las características. Con $C$, $I$ y $p$ fijos, el crecimiento respecto a $N$ es lineal. La predicción requiere $O(CMp)$, que es lineal en $M$ cuando $C$ y $p$ son fijos. La versión de referencia utilizó \texttt{OneVsRestClassifier} y \texttt{LogisticRegression(solver=lbfgs, max\_iter=2000, random\_state=42)}. L-BFGS puede detenerse antes del máximo al alcanzar el criterio de convergencia; su costo se expresa aproximadamente como $O(CI_{\mathrm{LBFGS}}Np)$ para entrenamiento y $O(CMp)$ para predicción (Scikit-learn Developers, 2026a).

\subsection{Evaluación y medición}
El desempeño se evaluó con Accuracy, F1 macro y Recall por clase. F1 macro promedia el F1 de cada clase sin ponderar por su frecuencia; Recall mide la proporción de casos reales de una categoría identificados correctamente (Pedregosa et al., 2011). Para los tiempos se midieron por separado \textit{fit} y \textit{predict} con \texttt{time.perf\_counter()}, después de una corrida de calentamiento y mediante siete repeticiones por configuración. Se reportan medianas en milisegundos y se ajustaron curvas para calcular $R^2$. El ajuste empírico describe los datos medidos, pero no demuestra por sí solo la complejidad asintótica; los tiempos también dependen del hardware, bibliotecas y carga del sistema.

\section{Resultados}
\subsection{Desempeño de clasificación}
La Tabla~\ref{tab:clasificacion} muestra que KNN obtuvo la mayor Accuracy y F1 macro. Las diferencias entre implementaciones propias y de \texttt{scikit-learn} fueron pequeñas.
\begin{table}[H]
\centering
\caption{Métricas globales de clasificación.}
\label{tab:clasificacion}
\scriptsize
\begin{tabular}{lcc}
\toprule
\textbf{Modelo} & \textbf{Accuracy} & \textbf{F1 macro}\\
\midrule
KNN propio & 90.12\% & 81.64\%\\
KNN scikit-learn & 89.88\% & 80.77\%\\
OvR propio & 87.76\% & 76.72\%\\
OvR scikit-learn & 87.29\% & 76.59\%\\
\bottomrule
\end{tabular}
\end{table}

\begin{table}[H]
\centering
\caption{Recall por clase.}
\label{tab:recall}
\scriptsize
\begin{tabular}{lccc}
\toprule
\textbf{Modelo} & \textbf{Normal} & \textbf{Sospechoso} & \textbf{Patológico}\\
\midrule
KNN propio & 97.28\% & 61.02\% & 71.43\%\\
KNN scikit-learn & 97.28\% & 61.02\% & 68.57\%\\
OvR propio & 95.77\% & 52.54\% & 71.43\%\\
OvR scikit-learn & 94.86\% & 52.54\% & 74.29\%\\
\bottomrule
\end{tabular}
\end{table}

La clase sospechosa presentó el Recall más bajo en todos los modelos. KNN propio obtuvo el mejor equilibrio global; la regresión logística de \texttt{scikit-learn} alcanzó el mayor Recall patológico (74.29\%).

\subsection{Tiempos de entrenamiento}
La Tabla~\ref{tab:fit} presenta las medianas para tres tamaños representativos. KNN propio tiene un ajuste casi constante porque almacena referencias. En cambio, el ajuste de KNN en \texttt{scikit-learn} fue mayor en tiempo medido (0.70--0.84 ms), debido a la validación y preparación de entradas. Para OvR, \texttt{scikit-learn} sí presentó menores tiempos de entrenamiento que la implementación propia.
\begin{table}[H]
\centering
\caption{Mediana del tiempo de entrenamiento (ms).}
\label{tab:fit}
\scriptsize
\begin{tabular}{crrrr}
\toprule
\textbf{Fracción} & \textbf{KNN P.} & \textbf{KNN SK} & \textbf{OvR P.} & \textbf{OvR SK}\\
\midrule
10\% & 0.0010 & 0.6999 & 115.8238 & 22.4092\\
50\% & 0.0010 & 0.7603 & 203.4343 & 29.9500\\
100\% & 0.0010 & 0.8375 & 320.7521 & 42.7461\\
\bottomrule
\end{tabular}
\end{table}

\begin{figure}[H]
\centering
\begin{subfigure}[t]{\columnwidth}
  \centering\grafica{KNN_entrenamiento.png}
  \caption{KNN: propio y scikit-learn.}
\end{subfigure}\vspace{1pt}
\begin{subfigure}[t]{\columnwidth}
  \centering\grafica{RegresionLogisticaOvR_entrenamiento.png}
  \caption{OvR: propio y scikit-learn.}
\end{subfigure}
\caption{Tiempos de entrenamiento según la fracción de datos utilizada.}
\label{fig:tiempos-fit}
\end{figure}

La Figura~\ref{fig:tiempos-fit} permite comparar por separado el crecimiento del tiempo de ajuste de KNN y OvR.

\subsection{Tiempos de predicción}
La Tabla~\ref{tab:predict} muestra que el tiempo de KNN propio aumenta notablemente con el tamaño del conjunto de entrenamiento. \texttt{scikit-learn} reduce el tiempo medido de predicción de KNN. Las dos versiones de OvR mantienen tiempos bajos en comparación con KNN propio.
\begin{table}[H]
\centering
\caption{Mediana del tiempo de predicción (ms).}
\label{tab:predict}
\scriptsize
\begin{tabular}{crrrr}
\toprule
\textbf{Fracción} & \textbf{KNN P.} & \textbf{KNN SK} & \textbf{OvR P.} & \textbf{OvR SK}\\
\midrule
10\% & 31.9714 & 3.9090 & 0.0542 & 0.5341\\
50\% & 692.6658 & 7.2958 & 0.0653 & 0.5137\\
100\% & 2926.5604 & 3.7136 & 0.0819 & 0.6229\\
\bottomrule
\end{tabular}
\end{table}

\begin{figure}[H]
\centering
\begin{subfigure}[t]{\columnwidth}
  \centering\grafica{KNN_prediccion.png}
  \caption{KNN: propio y scikit-learn.}
\end{subfigure}\vspace{1pt}
\begin{subfigure}[t]{\columnwidth}
  \centering\grafica{RegresionLogisticaOvR_prediccion.png}
  \caption{OvR: propio y scikit-learn.}
\end{subfigure}
\caption{Tiempos de predicción según la fracción de datos utilizada.}
\label{fig:tiempos-predict}
\end{figure}

La Figura~\ref{fig:tiempos-predict} destaca la diferencia entre KNN propio y las implementaciones de referencia al aumentar la muestra.

\subsection{Ajuste de las curvas temporales}
\begin{table}[H]
\centering
\caption{Coeficiente de determinación $R^2$ de los ajustes temporales.}
\label{tab:r2}
\scriptsize
\begin{tabular}{lcc}
\toprule
\textbf{Modelo} & \textbf{Fit} & \textbf{Predict}\\
\midrule
KNN propio & 0.0000 & 0.9982\\
KNN scikit-learn & 0.8189 & 0.0445\\
OvR propio & 0.9729 & 0.8926\\
OvR scikit-learn & 0.9070 & 0.1453\\
\bottomrule
\end{tabular}
\end{table}

Los valores de $R^2$ permiten evaluar qué tan bien las curvas seleccionadas representan la variación de los tiempos medidos dentro del intervalo experimental. Los valores altos indican un buen ajuste de la curva a los datos observados, pero no demuestran por sí mismos que la complejidad asintótica teórica del algoritmo sea correcta. Por otra parte, los valores bajos de $R^2$ obtenidos en las predicciones de KNN y OvR de scikit-learn indican que las curvas utilizadas no representan adecuadamente la variación de los tiempos en el intervalo medido. En el caso de KNN propio durante el entrenamiento, el valor $R^2 = 0$ se explica porque los tiempos registrados son prácticamente constantes, debido a que esta etapa únicamente almacena las referencias a los datos de entrenamiento y no realiza un procesamiento significativo. Por tanto, la complejidad teórica de cada algoritmo debe justificarse a partir de las operaciones que realiza y no únicamente mediante el ajuste estadístico de los tiempos experimentales.

\section{Discusión}
KNN obtuvo el mejor desempeño global, aunque el Recall de la clase sospechosa fue menor en todos los modelos. En tiempo, KNN propio combina un ajuste casi constante con una predicción costosa por el cálculo y ordenamiento de distancias; la biblioteca optimizada redujo considerablemente el tiempo de predicción. En OvR, la versión propia requirió más tiempo de entrenamiento que L-BFGS, mientras que ambas versiones tuvieron tiempos bajos de predicción. Los resultados experimentales son compatibles con el análisis de operaciones, pero no sustituyen el análisis teórico. El sistema clasifica categorías NSP y no predice directamente mortalidad fetal ni reemplaza la valoración clínica; su uso como apoyo requeriría validación externa y clínica.

\section{Conclusiones}
KNN presentó el mejor desempeño global de clasificación. Su implementación propia tuvo un entrenamiento prácticamente constante, pero una predicción cuyo tiempo creció con el tamaño de la muestra. \texttt{scikit-learn} fue notablemente más rápido en predicción de KNN y en entrenamiento de OvR. La clase sospechosa continúa siendo el principal reto de clasificación. Estos resultados se limitan al conjunto y al entorno experimental utilizados.

\balance
\section{Referencias}
\begingroup
\scriptsize
\setlength{\parindent}{0pt}
\setlength{\parskip}{1pt}
\noindent\hangindent=1.1cm\hangafter=1
Campos, D., y Bernardes, J. (2000). \textit{Cardiotocography} [Conjunto de datos]. UCI Machine Learning Repository. \url{https://doi.org/10.24432/C51S4N}

\noindent\hangindent=1.1cm\hangafter=1
NumPy Developers. (2026). \textit{numpy.argsort} (Versión 2.5) [Documentación de software]. \url{https://numpy.org/doc/stable/reference/generated/numpy.argsort.html}

\noindent\hangindent=1.1cm\hangafter=1
Pedregosa, F., Varoquaux, G., Gramfort, A., Michel, V., Thirion, B., Grisel, O., Blondel, M., Prettenhofer, P., Weiss, R., Dubourg, V., Vanderplas, J., Passos, A., Cournapeau, D., Brucher, M., Perrot, M., y Duchesnay, É. (2011). Scikit-learn: Machine learning in Python. \textit{Journal of Machine Learning Research, 12}, 2825--2830. \url{https://jmlr.org}

\noindent\hangindent=1.1cm\hangafter=1
Scikit-learn Developers. (2026a). \textit{LogisticRegression} (Versión 1.9.1) [Documentación de software]. \url{https://scikit-learn.org/stable/modules/generated/sklearn.linear_model.LogisticRegression.html}

\noindent\hangindent=1.1cm\hangafter=1
Scikit-learn Developers. (2026b). \textit{Nearest Neighbors} (Versión 1.9.1) [Documentación de software]. \url{https://scikit-learn.org/stable/modules/neighbors.html}

\endgroup

\end{document}
